\documentclass[11pt]{article}
\usepackage{preamble}
\renewcommand{\answer}[1]{}

\begin{document}

\ladderheading{AP Statistics \textperiodcentered\ Unit 2 \textperiodcentered\ Topic 2.9 \textperiodcentered\ B: 2026-11-02 \textperiodcentered\ E: 2026-11-23}{Same Mean, Different Deadline Risk}

\focusbanner{TODAY'S PROBLEM}

Suppose a dispatcher models delivery time using eight equally likely outcomes for each route, shown in the dotplots. An analyst reports that both routes have mean 30 minutes, with $\sigma_A=3$ minutes and $\sigma_B=12$ minutes.\par\smallskip The dispatcher says, \emph{The routes have the same mean, so it does not matter which one a customer gets.} A customer has 40 minutes before an appointment and places a high cost on being late.\par
\begin{center}
\begin{tikzpicture}
\begin{axis}[
  width=0.90\linewidth,
  height=0.68in,
  xmin=5,
  xmax=55,
  xtick={10,15,20,25,30,35,40,45,50},
  ymin=0,
  ymax=3,
  ytick=\empty,
  axis y line=none,
  axis x line=bottom,
  title={\textbf{Route A}},
  xlabel={Delivery time (minutes)},
  tick label style={font=\scriptsize},
  label style={font=\scriptsize},
  title style={font=\small}
]
\addplot[only marks,mark=*,mark size=2.4pt,color=dayblue] coordinates
  {(25,1) (27,1) (29,1) (29,2) (31,1) (31,2) (33,1) (35,1)};
\end{axis}
\end{tikzpicture}\par\vspace{-0.08in}
\begin{tikzpicture}
\begin{axis}[
  width=0.90\linewidth,
  height=0.68in,
  xmin=5,
  xmax=55,
  xtick={10,15,20,25,30,35,40,45,50},
  ymin=0,
  ymax=3,
  ytick=\empty,
  axis y line=none,
  axis x line=bottom,
  title={\textbf{Route B}},
  xlabel={Delivery time (minutes)},
  tick label style={font=\scriptsize},
  label style={font=\scriptsize},
  title style={font=\small}
]
\addplot[only marks,mark=*,mark size=2.4pt,color=dayblue] coordinates
  {(10,1) (18,1) (26,1) (26,2) (34,1) (34,2) (42,1) (50,1)};
\end{axis}
\end{tikzpicture}
\end{center}
\begin{firsttakebox}{FIRST TAKE --- before the video (5 min)}
Which route would you assign to this customer? Cite one number or dotplot feature.
\workspace{0.9in}
\end{firsttakebox}

\begin{callout}{\IconBook\ MEAN, STANDARD DEVIATION, AND SHIFTS (keep on the page)}{calloutgreen}
For a discrete random variable, $\mu_X=\sum xP(X=x)$ and
$\sigma_X=\sqrt{\sum(x-\mu_X)^2P(X=x)}$. Interpret $\mu_X$ as the long-run average and
$\sigma_X$ as a typical distance from that mean, in context and with units. If $Y=X+c$, then
$\mu_Y=\mu_X+c$ while $\sigma_Y=\sigma_X$: adding the same constant shifts every outcome but does not change spread.
Match the parameter to the decision; equal centers do not imply equal tail risk.
\end{callout}

\newpage
\textbf{Finish after the video:}\par\smallskip

\dokbadge{1}~\textbf{(a)} For each route, state its minimum and maximum modeled time and count the outcomes longer than 40 minutes.\par
\workspace{0.7in}

\dokbadge{2}~\textbf{(b)} Verify the reported mean and population standard deviation for each route from the eight equally likely outcomes. Show the sums of squared deviations.\par
\workspace{1.4in}

\dokbadge{3}~\textbf{(c)} Decide which route better matches this customer's priority. Cite the parameter and dotplot feature that settle your choice, and explain what the dispatcher's mean-only argument would mislead a reader into believing. Then suppose every delivery gets 5 minutes of loading time: state both transformed means and standard deviations and decide whether your route choice changes.\par
\workspace{2.0in}

\begin{sentenceframebox}\raggedright\textbf{Frame:} Route \blank[0.4in] fits the customer because $\sigma=\blank[0.5in]$ and the dotplot shows \blank[2.2in].\end{sentenceframebox}

\begin{sentenceframebox}\raggedright\textbf{Frame:} Adding 5 changes each mean to \blank[0.5in] but leaves the SDs at \blank[0.5in] and \blank[0.5in] because \blank[2.0in].\end{sentenceframebox}

\begin{summaryexitbox}{EXIT --- turn this in}
One thing the video changed about my first take: \hrulefill\par
\workspace{0.4in}
\end{summaryexitbox}

\end{document}
